\makeabstract
    {
    Chemical Comparison of Two Compact Cores in the High-Mass Star-Forming Region IRAS 07299-1651
    }
    {
    Damien Algernon\textsuperscript{1}, Dr. M\'elisse Bonfand\textsuperscript{2}, Dr. Robin Garrod\textsuperscript{2}, Dr. Mark Marshall\textsuperscript{1}
    }
    {
    \textsuperscript{1} Department of Chemistry, Amherst College, Amherst, MA\\
\textsuperscript{2} Department of Chemistry, University of Virginia, Charlottesville, VA
    }
    {
    High-mass star-forming regions contain warm, dense gas with rich molecular emission, but neighboring protostellar cores may experience different chemical and physical conditions. This study compares two compact continuum cores in the Infrared Astronomical Satellite (IRAS) source 07299-1651 to determine whether their molecular inventories and kinematic properties are similar or chemically differentiated.
    }
    {
    Submillimeter Array (SMA) continuum maps and spectral cubes were analyzed across 24 spectral windows. Astrodendro was used to identify the two continuum peaks, and their pixel positions were used to extract spectra with Common Astronomy Software Applications (CASA). Emission features were counted, and representative profiles were compared using one- and two-component Gaussian models. The eXtended CASA Line Analysis Software Suite (XCLASS) generated synthetic spectra and fitted molecular column densities, velocities, and linewidths for each core.
    }
    {
    XCLASS assigned molecular carriers to 96\% and 94\% of the counted emission features in Cores 1 and 2, respectively. The cores display broadly similar molecular inventories, but their fitted column densities vary by species. Several molecules have higher fitted column densities toward Core 2, whereas sulfur monoxide and the cyano radical are enhanced toward Core 1. Mean fitted velocities of 17.1 km/s for Core 1 and 16.8 km/s for Core 2 are close to the reported systemic velocity of 16.6 km/s. Core 2 has a broader mean molecular linewidth than Core 1, 6.1 versus 3.6 km/s. A two-component Gaussian model better reproduces a representative line. The model is favored over a single-component model, indicating that blended or unresolved velocity components may contribute to the observed profiles.
    }
    {
    The two compact cores in IRAS 07299-1651 are chemically related but not identical. Species-dependent column densities and the broader linewidth of Core 2 indicate differences in their local chemical or kinematic environments, while their mean velocities remain consistent with the surrounding region. Future work will refine uncertain molecular fits and apply the analysis pipeline to nine additional high- and intermediate-mass star-forming regions observed with the SMA.
    }

    \newpage
    